% Documentación de la Etapa 1 - generado a partir de DOCUMENTACION.md
% Compilar con:  pdflatex DOCUMENTACION.tex
\documentclass[11pt,a4paper]{article}

\usepackage[utf8]{inputenc}
\usepackage[margin=2.5cm]{geometry}
\usepackage{url}

\title{Etapa 1 - Analizador Léxico de MiniJava}
\author{Compiladores e Intérpretes --- UNS, 2026}
\date{}

% Esta instalacion TeX no tiene las fuentes TS1 (text companion): se evitan los
% simbolos que las requieren (vinetas de itemize, llaves via \{ en \texttt).
\renewcommand{\labelitemi}{$\bullet$}

\begin{document}

\sloppy

\maketitle

\noindent
Implementación de un analizador lexico para MiniJava junto con un módulo principal de
consola. El analizador se construye a partir de un autómata finito determinista
implementado manualmente (un método por estado).

\section{Estructura del proyecto}

{\footnotesize
\begin{verbatim}
codigo/
|-- Compilar.sh / Compilar.bat          Compila y genera Compilador.jar
|-- CorrerTests.sh / CorrerTests.bat    Compila y ejecuta los testers JUnit
|-- lib/                                Jars de JUnit 4 y Hamcrest (se autodescargan)
|-- src/main/java/
|   |-- moduloprincipal/
|   |   \-- ModuloPrincipal.java        Modulo principal (interfaz con el usuario)
|   |-- analizadorlexico/
|   |   |-- AnalizadorLexico.java       Automata finito determinista
|   |   |-- Token.java                  Representacion de un token (nombre, lexema, linea)
|   |   \-- ExcepcionLexica.java        Excepcion propagada ante un error lexico
|   \-- sourcemanager/
|       |-- SourceManager.java          Interfaz provista por la catedra
|       |-- SourceManagerImpl.java      Implementacion provista (por lineas, referencia)
|       \-- SourceManagerEficiente.java Implementacion propia por caracteres (se usa)
|-- src/test/java/test/                 Testers JUnit provistos por la catedra
\-- resources/
    |-- sinErrores/                     Casos de prueba sin errores
    \-- conErrores/                     Casos de prueba con errores
\end{verbatim}
}

\section{Componentes}

\subsection{Módulo principal (ModuloPrincipal)}

Recibe por parámetro la ruta del archivo fuente MiniJava (acepta cualquier extensión):

\begin{verbatim}
java -jar Compilador.jar programa1.java
\end{verbatim}

Asocia el archivo a un \texttt{SourceManagerEficiente}, crea el analizador léxico y le solicita
tokens hasta recibir el token de fin de archivo. Cada token se imprime por pantalla con el
formato \texttt{(NombreToken,Lexema,NroLinea)}. Al finalizar sin errores se imprime
\texttt{[SinErrores]}.

Ante un error léxico, la excepción \texttt{ExcepcionLexica} es capturada por el módulo principal,
que reporta el mensaje de error y el código de error, finalizando la ejecución.

\subsection{Manejador de archivos eficiente (SourceManagerEficiente)}

Implementa la interfaz \texttt{SourceManager} leyendo el fuente de a caracteres (en lugar de
líneas como la implementación provista). Características:

\begin{itemize}
  \item Lectura por caracteres con \texttt{InputStreamReader} (UTF-8) sobre \texttt{PushbackReader}.
  \item Reconoce las tres convenciones de salto de línea (LF, CR y CRLF), normalizándolas a un
        único \verb!\n!. La secuencia \verb!\r\n! no se
        contabiliza como dos saltos.
  \item \texttt{getLineNumber()} devuelve la línea del último caracter devuelto. Un caracter de salto
        de línea pertenece a la línea que termina, y el número de línea se incrementa al leer
        el siguiente caracter. De esta forma el token EOF reporta la línea del final físico
        del archivo (por ejemplo, un archivo de 21 líneas que termina con salto de línea
        reporta el EOF en la línea 22).
\end{itemize}

La implementación provista por la cátedra (\texttt{SourceManagerImpl}) se incluye como
referencia, pero no se utiliza en la ejecución porque su conteo de líneas difiere en el
token EOF cuando el archivo termina con un salto de línea.

\subsection{Analizador lexico (AnalizadorLexico)}

Autómata finito determinista implementado con un método por estado. El estado inicial
\texttt{e0} descarta espacios en blanco y comentarios, y despacha al estado correspondiente
según el primer caracter del lexema. Cada método consulta \texttt{caracterActual} (un caracter
de adelanto que nunca se consume al finalizar un token), lo agrega al lexema y avanza
leyendo el próximo caracter.

Las palabras reservadas se resuelven consultando una tabla: al construir un identificador
de método o variable se verifica si su lexema corresponde a una palabra reservada y, en
ese caso, se devuelve el token de palabra reservada correspondiente.

\section{Tokens del lenguaje}

Se adoptaron los siguientes nombres para los tokens.

\subsection{Palabras reservadas}

Se reconocen con un identificador de método o variable y se devuelven con nombre
\verb!pr_<palabra>!:

\noindent
\texttt{class} \texttt{extends} \texttt{interface} \texttt{implements} \texttt{static}
\texttt{boolean} \texttt{char} \texttt{int} \texttt{void}
\texttt{public} \texttt{if} \texttt{else} \texttt{while} \texttt{return} \texttt{var}
\texttt{this} \texttt{new} \texttt{null} \texttt{true} \texttt{false}

\medskip
\noindent
Los nombres de tokens siguen el ejemplo de las pautas de la Etapa 1: \texttt{litString}
para los literales string, \texttt{idMV} para identificadores de método o variable,
\verb!pr_<palabra>! para palabras reservadas, \texttt{op<simbolo>} para operadores y
\texttt{llaveC} para puntuación.

\subsection{Identificadores}

\begin{center}
\begin{tabular}{|l|l|l|}
\hline
\textbf{Token} & \textbf{Expresión regular} & \textbf{Descripción} \\ \hline
\texttt{idClase} & \verb![A-Z][A-Za-z0-9_]+! & Clase o interfaz \\ \hline
\texttt{idGen} & \texttt{[A-Z]} & Parámetro de tipo genérico (una sola mayúscula) \\ \hline
\texttt{idMV} & \verb![a-z][A-Za-z0-9_]*! & Método o variable \\ \hline
\end{tabular}
\end{center}

\subsection{Literales}

\begin{center}
{\footnotesize
\setlength{\tabcolsep}{4pt}
\begin{tabular}{|l|l|p{5.5cm}|}
\hline
\textbf{Token} & \textbf{Expresión regular} & \textbf{Restricciones} \\ \hline
\texttt{litInt} & \verb![0-9]{1,9}! & Máximo 9 dígitos; más de 9 es error \\ \hline
\texttt{litChar} & \verb!'[^\\'\n]'! o \verb!'\\[^\n]'! &
  Un caracter normal, o barra invertida seguida de cualquier caracter \\ \hline
\texttt{litString} & \verb!"([^"\\\n]|\\[^\n])*"! &
  Sin saltos de línea; \verb!\"! no cierra el literal \\ \hline
\end{tabular}
}
\end{center}

Los literales \texttt{true}, \texttt{false} y \texttt{null} se representan con las palabras reservadas
\verb!pr_true!, \verb!pr_false! y \verb!pr_null!.

\subsection{Operadores}

\noindent
\texttt{op+} \texttt{op-} \texttt{op*} \texttt{op/} \texttt{op\%} \texttt{op<} \texttt{op>}
\texttt{op!} \texttt{op=} \texttt{op==} \texttt{op>=} \texttt{op<=} \texttt{op!=}
\texttt{op\&\&} \verb!op||! \texttt{op++} \texttt{op--}

\medskip
\noindent
Los operadores se reconocen con el criterio del lexema más largo (por ejemplo, \texttt{>=}
antes que \texttt{>}). Los caracteres \texttt{\&} y \verb!|! no forman operadores por sí solos:
un \texttt{\&} o \verb!|! no seguido de otro igual es un error léxico.

\subsection{Puntuación}

\begin{center}
\begin{tabular}{|l|l|l|l|l|l|}
\hline
\texttt{parA} & \texttt{parC} & \texttt{llaveA} & \texttt{llaveC} & \texttt{corcheteA} & \texttt{corcheteC} \\
\texttt{(} & \texttt{)} & \verb!{! & \verb!}! & \texttt{[} & \texttt{]} \\ \hline
\texttt{puntoComa} & \texttt{coma} & \texttt{punto} & \texttt{dosPuntos} & & \\
\texttt{;} & \texttt{,} & \texttt{.} & \texttt{:} & & \\ \hline
\end{tabular}
\end{center}

\subsection{Fin de archivo}

Token \texttt{EOF} con lexema \verb!$!.

\section{Elementos ignorados}

\begin{itemize}
  \item Espacios en blanco: espacio, tabulador y salto de línea (LF, CR o CRLF).
  \item Comentarios de una línea: desde \texttt{//} hasta el final de la línea.
  \item Comentarios multilínea: desde \texttt{/*} hasta \texttt{*/}.
\end{itemize}

\section{Errores lexicos}

Al detectar un error se lanza \texttt{ExcepcionLexica} con el lexema armado hasta el momento y
el número de línea, y el módulo principal finaliza la ejecución reportando:

\begin{verbatim}
Error Lexico en linea <N>: <mensaje>
[Error:<lexema>|<N>]
\end{verbatim}

Para que el código de error quede en una sola línea, en el caso de un comentario
multilínea sin cerrar el lexema reportado es el que abre el comentario, \texttt{/*}.

Mensajes según la naturaleza del error:

\begin{center}
\begin{tabular}{|l|l|}
\hline
\textbf{Caso} & \textbf{Mensaje} \\ \hline
Caracter inválido (no es símbolo válido) & \texttt{<lexema> no es un simbolo valido} \\ \hline
String sin comilla de cierre & \texttt{literal string mal formado} \\ \hline
Caracter mal formado & \texttt{literal caracter mal formado} \\ \hline
Entero con más de 9 dígitos & \texttt{literal entero demasiado largo} \\ \hline
Comentario multilínea sin cerrar & \texttt{comentario sin cerrar} \\ \hline
\end{tabular}
\end{center}

\section{Formato de salida}

\subsection{Análisis exitoso}

Un token por línea con el formato \texttt{(NombreToken,Lexema,NroLinea)}. El lexema incluye los
delimitadores para strings y caracteres (por ejemplo, \texttt{(litString,"hola",2)} o
\verb!(litChar,'\n',3)!). Al final se imprime el token de fin de archivo y:

\begin{verbatim}
[SinErrores]
\end{verbatim}

\subsection{Análisis con errores}

\begin{verbatim}
Error Lexico en linea 2: # no es un simbolo valido
[Error:#|2]
\end{verbatim}

Toda la salida se realiza con \texttt{System.out}.

\section{Compilacion y ejecucion}

No se usa ninguna herramienta de build: solo se necesitan \texttt{javac}, \texttt{jar} y
\texttt{java} de cualquier JDK 21 o superior (probado con JDK 25).

{\footnotesize
\begin{verbatim}
./Compilar.sh             # compila y genera Compilador.jar  (Windows: Compilar.bat)
java -jar Compilador.jar resources/sinErrores/lexSinErrores01.java
./CorrerTests.sh          # compila y ejecuta los testers JUnit  (Windows: CorrerTests.bat)
\end{verbatim}
}

\texttt{CorrerTests.sh} descarga automáticamente \texttt{junit-4.13.2.jar} y
\texttt{hamcrest-core-1.3.jar} en \texttt{lib/} si no están (requiere \texttt{curl});
si no hay red, se pueden descargar a mano de Maven Central y dejarlos en esa carpeta.

\subsection{Comandos manuales equivalentes}

Compilar (desde \texttt{codigo/}):

{\footnotesize
\begin{verbatim}
javac -d build/classes \
  src/main/java/moduloprincipal/*.java \
  src/main/java/analizadorlexico/*.java \
  src/main/java/sourcemanager/*.java
jar cfe Compilador.jar moduloprincipal.ModuloPrincipal -C build/classes .
\end{verbatim}
}

Correr los tests (desde \texttt{codigo/}, con JUnit 4.13.2 y Hamcrest 1.3 en \texttt{lib/}):

{\footnotesize
\begin{verbatim}
javac -cp "build/classes:lib/junit-4.13.2.jar:lib/hamcrest-core-1.3.jar" \
  -d build/test-classes src/test/java/test/*.java

java -cp "build/classes:build/test-classes:lib/junit-4.13.2.jar:\
lib/hamcrest-core-1.3.jar" org.junit.runner.JUnitCore \
  test.TesterDeCasosSinErrores test.TesterDeCasosConErrores
\end{verbatim}
}

En Windows el separador del classpath es \texttt{;} en lugar de \texttt{:}.

Los testers (\texttt{TesterDeCasosSinErrores} y \texttt{TesterDeCasosConErrores}) usan el working
directory \texttt{codigo/} para acceder a \texttt{resources/sinErrores/} y
\texttt{resources/conErrores/} (los scripts ya se posicionan ahí; los comandos manuales
deben correrse desde \texttt{codigo/}).

\section{Tests provistos por la catedra}

Casos de prueba incluidos (copiados del material de la cátedra):

\begin{itemize}
  \item \texttt{resources/sinErrores/lexSinErrores01.java} - identificadores, strings, chars y EOF.
  \item \texttt{resources/sinErrores/lexSinErrores02.java} - ejemplo de la convención de salida.
  \item \texttt{resources/sinErrores/lexSinErrores03.java} - archivo solo con comentarios.
  \item \texttt{resources/conErrores/lexConErrores01.java} - caracter inválido \texttt{\#}.
  \item \texttt{resources/conErrores/lexConErrores02.java} - string sin comilla de cierre.
  \item \texttt{resources/conErrores/lexConErrores03.java} - múltiples errores en una corrida.
\end{itemize}

\section{Logros}

\begin{itemize}
  \item \textbf{Manejador de Archivos Eficiente}: \texttt{SourceManagerEficiente} lee de a
        caracteres y conforma con la interfaz \texttt{SourceManager}, sin modificar la
        interfaz provista.
  \item \textbf{Multi-detección de Errores Léxicos}: el módulo principal no finaliza ante
        el primer error léxico; captura la excepción, reporta el error y continúa pidiendo
        tokens hasta el \texttt{EOF}. El léxico consume el carácter ofensor en los estados
        de error para garantizar el progreso de la recuperación. Sólo si la corrida no
        registró errores se imprime \texttt{[SinErrores]}.
\end{itemize}

\end{document}
